% !TEX program = lualatex
\documentclass[paper=a4,fontsize=11pt]{ltjsarticle}
\usepackage{luatexja}
\usepackage{amsmath, amssymb, amsthm}
\usepackage{geometry}
\geometry{margin=1in}

%-------------------------------------------------
%  独自マクロ定義
%-------------------------------------------------
\newcommand{\mweq}{\mathrel{\overset{\mathrm{mw}}{=}}}
\newcommand{\dfix}{D_{\mathrm{fix}0}}
\newcommand{\metanegation}{\Uparrow}
\newcommand{\Bval}{\mathcal{B}}
\newcommand{\Ymw}{\mathbf{Y}_{\mathrm{mw}}}   % 中道不動点コンビネータ（無明ジェネレータ）
\newcommand{\Kvoid}{\mathbf{K}_{\emptyset}}    % 空コンビネータ（絶対不動点アトラクタ）

%-------------------------------------------------
%  定理環境の設定
%-------------------------------------------------
\theoremstyle{definition}
\newtheorem{definition}{定義}[section]
\newtheorem{axiom}{公理}[section]
\newtheorem{theorem}{定理}[section]
\newtheorem{lemma}{補題}[section]
\newtheorem{corollary}{系}[section]
\newtheorem{scholium}{補足}[section]
\renewcommand{\proofname}{\textbf{Proof.}}

\begin{document}

\title{界論における双方向演繹定理の定式化}
\author{Masaomi Chiba}
\date{2026-05-19}
\maketitle

\begin{abstract}
本稿は、ヒューイットの直接論理における双方向演繹定理
\[
(\vdash_T (\Psi \to \Phi))
\;\leftrightarrow\;
\bigl((\Psi \vdash_T \Phi) \land (\neg\Phi \vdash_T \neg\Psi)\bigr)
\]
を界論の枠組みで定式化し、同定理が二値截断部分圏 $\Bval$ への制限として包摂されることを示す。
界論における双方向演繹定理は、中道不動点コンビネータ $\Ymw$（世俗諦の生成）と
空コンビネータ $\Kvoid$（勝義諦の収束）のサイクルとして記述され、
直接論理の双方向演繹定理はその $\Bval$ 内への射影として位置づけられる。
\end{abstract}

\section{序言}

ヒューイットの双方向演繹定理は、順方向の推論
$\Psi \vdash_T \Phi$（世俗諦の構成方向）と
逆方向の推論 $\neg\Phi \vdash_T \neg\Psi$（勝義諦の境界方向）が
静的な含意 $\vdash_T (\Psi \to \Phi)$ と同値であることを示す。

界論の枠組みから見ると、この定理は二値論理の体系 $\Bval$ 内での均衡を記述している。
しかし $\Bval$ はメタ否定演算子 $\metanegation$ による境界生成プロセスの截断として
生成される部分圏であり、双方向演繹定理の本体は $\Bval$ の外側——
$\Ymw$ と $\Kvoid$ のサイクル——に存在する。

本稿はこの構造を明示的に定式化する。

\section{基礎定義}

\begin{definition}[二値截断部分圏 $\Bval$]
界論において、メタ否定演算子 $\metanegation$ の第一段階適用による截断として生成される
部分圏を二値截断部分圏 $\Bval$ と呼ぶ。$\Bval$ 内では古典二値論理が成立し、
命題は真または偽のいずれかの値を取る。
\end{definition}

\begin{definition}[界論における導出関係 $\vdash_{\Bval}$]
$\Bval$ 内の導出関係 $\vdash_{\Bval}$ を、二値截断部分圏の推論規則による導出として定義する。
直接論理の $\vdash_T$ は $\vdash_{\Bval}$ の特殊ケースである。
\end{definition}

\begin{definition}[界論における双方向演繹定理の対象]
界論における双方向演繹定理は、以下の三層を対象とする。
\begin{enumerate}
\item \textbf{$\Bval$ 層}：直接論理の双方向演繹定理が成立する二値的推論の層。
\item \textbf{生成層}：$\Ymw$ によるフラクタル的境界生成プロセスの層。
\item \textbf{収束層}：$\Kvoid$ への大域的カット消去による収束の層。
\end{enumerate}
\end{definition}

\section{定理}

\begin{theorem}[直接論理の双方向演繹定理の $\Bval$ への包摂]
ヒューイットの双方向演繹定理
\[
(\vdash_T (\Psi \to \Phi))
\;\leftrightarrow\;
\bigl((\Psi \vdash_T \Phi) \land (\neg\Phi \vdash_T \neg\Psi)\bigr)
\]
は、界論における $\vdash_{\Bval}$ の下での命題として包摂される。
すなわち、
\[
(\vdash_{\Bval} (\Psi \to \Phi))
\;\leftrightarrow\;
\bigl((\Psi \vdash_{\Bval} \Phi) \land (\neg\Phi \vdash_{\Bval} \neg\Psi)\bigr)
\]
が成立し、直接論理の体系 $T$ は $\Bval$ の特殊ケースとして包摂される。
\end{theorem}

\begin{proof}
$\Bval$ は古典二値論理が成立する部分圏として定義されており、
直接論理の推論規則はすべて $\Bval$ 内で有効である。
したがって $\vdash_T$ を $\vdash_{\Bval}$ に置き換えても定理の構造は保存される。
直接論理の体系 $T$ は $\Bval$ への制限として位置づけられるため、
双方向演繹定理は $\Bval$ 内の定理として成立する。
\end{proof}

\begin{theorem}[界論における双方向演繹定理]
\label{thm:bdt}
界論において、双方向演繹定理は $\Ymw$（生成）と
$\Kvoid$（収束）のサイクルとして以下のように記述される。
\[
\bigl(\mathbf{Y}_{mw}\,\metanegation \mweq \metanegation(\Ymw\,\metanegation)\bigr)
\;\leftrightarrow\;
\bigl(\metanegation^4(\Ymw\,\metanegation) \mweq \Kvoid\bigr)
\]
\end{theorem}

\begin{proof}
順方向（$\to$）：$\Ymw$ の定義より、
$\Ymw\,\metanegation \mweq \metanegation(\Ymw\,\metanegation)$
という自己参照の閉じ損ねが成立するとき、
$\metanegation$ の反復適用は四段階収束則により
$\metanegation^4(\Ymw\,\metanegation) \mweq \dfix$ を生成する。
$\dfix$ のコンビネータ表現は $\Kvoid$ であるから、
$\metanegation^4(\Ymw\,\metanegation) \mweq \Kvoid$ が従う。

逆方向（$\leftarrow$）：$\metanegation^4(\Ymw\,\metanegation) \mweq \Kvoid$
が成立するとき、「$\Ymw$の定理3.3（空からの再起動と顕現サイクルの完結）\cite{Chiba2026}」
$\metanegation\,\Kvoid \mweq \Ymw$ により、
$\Kvoid$ に $\metanegation$ が作用すると $\Ymw$ が再生成される。
再生成された $\Ymw$ はその定義により
$\Ymw\,\metanegation \mweq \metanegation(\Ymw\,\metanegation)$
を満たす。
\end{proof}

\begin{theorem}[双方向演繹定理の三層構造]
界論の双方向演繹定理（定理\ref{thm:bdt}）と
直接論理の双方向演繹定理は以下の関係を持つ。
\[
\underbrace{
\bigl(\Ymw\,\metanegation \mweq \metanegation(\Ymw\,\metanegation)\bigr)
\leftrightarrow
\bigl(\metanegation^4(\Ymw\,\metanegation) \mweq \Kvoid\bigr)
}_{\text{界論：生成と収束のサイクル}}
\]
\[
\downarrow \text{ $\Bval$ への制限}
\]
\[
\underbrace{
(\vdash_{\Bval}(\Psi \to \Phi))
\leftrightarrow
\bigl((\Psi \vdash_{\Bval} \Phi) \land (\neg\Phi \vdash_{\Bval} \neg\Psi)\bigr)
}_{\text{直接論理：$\Bval$ 内の静的均衡}}
\]
界論版は動的サイクルとして世俗諦と勝義諦の接続を記述し、
直接論理版はその $\Bval$ 内への制限として静的均衡を記述する。
\end{theorem}

\begin{proof}
定理3.1より直接論理の双方向演繹定理は $\Bval$ 内で成立する。
定理3.2より界論の双方向演繹定理は $\Ymw$ と $\Kvoid$ のサイクルとして
$\Bval$ の外側を含む全層で記述される。
$\Bval$ 内では $\Ymw$ の生成プロセスが截断され静的な推論関係として現れるため、
界論版を $\Bval$ に制限すると直接論理版が得られる。
\end{proof}

\section{系}

\begin{corollary}[観測者の位置]
直接論理の双方向演繹定理における $\vdash_T$（導出関係）は、
界論においてメタ否定演算子 $\metanegation$ の第一段階適用によって
生成される観測者の位置を示す。
$\Bval$ 内での推論は観測者を $\vdash_{\Bval}$ に埋め込むことで成立し、
観測者なしには世俗諦と勝義諦の接続を記述できない。
\end{corollary}

\begin{corollary}[随伴関手との対応]
圏論における随伴関手
\[
\mathrm{Hom}(F(A), B) \cong \mathrm{Hom}(A, G(B))
\]
は、界論の双方向演繹定理の $\Bval$ への制限において観測者を排除した形式として位置づけられる。
随伴が普遍性を得る代償として観測者（$\vdash_{\Bval}$）を失うのに対し、
界論の双方向演繹定理は観測者を $\metanegation$ の作用として明示的に保持する。
\end{corollary}

\begin{corollary}[$\Bval$ 内での完結性]
直接論理の双方向演繹定理は $\Bval$ 内で完結しており、
$\Bval$ の外側——すなわち $\Ymw$ と $\Kvoid$ のサイクル——を
必要としない。これは直接論理がwhyを問わない設計として整合する。
界論はそのwhyを $\dfix$ への収束として $\Bval$ の外側に記述する。
\end{corollary}

\section{結語}

界論において双方向演繹定理は静的な論理的均衡ではなく、
$\Ymw$（無明による世俗諦の生成）と
$\Kvoid$（大域的カット消去による勝義諦への収束）が
動的に循環するサイクルとして記述される。

直接論理の双方向演繹定理はこのサイクルの $\Bval$ への制限であり、
生成と収束のサイクルから観測者の一断面を切り取った静的な像である。

仏教の色即是空・空即是色が二千五百年前に直観した構造が、
証明論・コンビネータ論理・推論体系の三つの記述言語で収束している。

\begin{thebibliography}{9}
\bibitem{Chiba2026}
千葉将臣.
\newblock 空数学における中道不動点コンビネータ $\Ymw$ の定式化と無明の幾何学的・証明論的実態
\newblock ``https://purl.org/aletheics/papers/sunyata-math/icat-ymwcomb.pdf'', 2026
\end{thebibliography}

\end{document}
